\title{Lecture 12\\User interfaces of intelligent systems \vspace{-2em}}
%\author[]{Shunkevich D.V.}
%\institute[]{Belarusian State University of Informatics and Radioelectronics}

\begin{frame}
	\titlepage
\end{frame}

\begin{frame}{Lecture Contents}
	\topline
	\justifying
	
	
	User interfaces as intelligent systems. External languages of user communication: universal and specialized. Editors of external information structures.
	
\end{frame}

\begin{frame}{Interface as a subsystem\\ of a cybernetic system}
	\topline
	\justifying
	
	The interface of a cybernetic system is considered as a special subsystem that ensures the interaction of the system with the outside world: the user, other systems, technical devices and the physical environment.
	
	\begin{textitemize}
		\item The interface is not limited to on-screen dialogue with a person.
		\item It provides for the reception, interpretation, transmission and presentation of external influences and messages.
		\item In an intelligent system, the interface must be consistent with the models of knowledge, context, and goals of the activity.
	\end{textitemize}
\end{frame}

\begin{frame}{External environment for a cybernetic system}
	\topline
	\justifying
	
	
	The external environment includes not only the human operator, but also many other sources and recipients of interaction.
	
	\begin{textitemize}
		\item Users and user groups.
		\item Other cybernetic and information systems.
		\item Sensors, actuators, robotic components.
		\item Documents, databases, signals, images, scenes of the outside world.
		\item Communication channels and exchange protocols.
	\end{textitemize}
\end{frame}

\begin{frame}{Why is the interface especially important for intelligent systems?}
	\topline
	\justifying
	
	
	For conventional programs, an interface is often understood as a means of input and output. For intelligent systems, this is insufficient, as interaction must be meaningfully interpretable, context-sensitive, and adaptive.
	
	\begin{textitemize}
		\item It is necessary to correlate messages with domain models.
		\item It is necessary to support uncertainty, incompleteness and ambiguity.
		\item It is necessary to take into account the user's state, tasks and interaction situation.
		\item It is important to provide an explanation for the system's decisions.
	\end{textitemize}
\end{frame}

\begin{frame}{User Interface and Intelligent Interface}
	\topline
	\justifying
	
	These concepts are close, but not identical.
	
	\begin{textitemize}
		\item The user interface of an intelligent system is a subsystem of interaction between the system and the user.
		\item An intelligent interface is a broader concept; it uses knowledge, context models, interpretation and adaptation mechanisms.
		\item Intelligence can be manifested in both human-machine and intersystem and sensory-executive interactions.
		\item Therefore, not every interface of an intelligent system is intelligent in the strong sense, but advanced intelligent systems usually require precisely such interfaces.
	\end{textitemize}
\end{frame}

\begin{frame}{Working Definition of an Intelligent Interface}
	\topline
	\justifying
	
	
	An intelligent interface is an interface subsystem of a cybernetic system that, based on knowledge of the subject domain, user, context and goals of the activity, interprets input messages, generates output messages and adapts the interaction process.
	
	\begin{textitemize}
		\item Uses knowledge models, not just a set of commands and forms.
		\item Supports semantic interpretation of messages.
		\item Capable of personalization, explanation and adaptation.
		\item Can work with various modalities: text, speech, graphics, gestures, sensory signals.
	\end{textitemize}
\end{frame}

\begin{frame}{Working Definition of the User Interface of an Intelligent System}
	\topline
	\justifying
	
	
	The user interface of an intelligent system is a part of the interface subsystem that directly organizes two-way interaction between the system and a person at the level of information presentation, input, interpretation of messages, support of actions and feedback.
	
	\begin{textitemize}
		\item It should ensure not only convenience, but also substantive accuracy of communication.
		\item It focuses on message semantics, user intent modeling, and explainability.
		\item It may include graphic, text, speech, multimodal and mixed means.
	\end{textitemize}
\end{frame}

\begin{frame}{Interface Subsystem Boundaries}
	\topline
	\justifying
	
	
	The interface subsystem of an intelligent system is located between the internal problem-solving mechanisms and the outside world.
	
	\begin{textitemize}
		\item As input, it receives external actions, messages and observations.
		\item In the inner layer, it relates them to knowledge, context, scenarios and rules.
		\item At the output, it generates external information structures that are understandable to the user or another system.
		\item Often includes monitoring, dialogue, explanation, and attention management tools.
	\end{textitemize}
\end{frame}

\begin{frame}{Three views of the interface}
	\topline
	\justifying
	
	
	It is useful to consider the same interface from three interrelated perspectives.
	
	\begin{textitemize}
		\item Functional: what classes of interaction problems does it solve.
		\item Informational-semiotic: what external languages and constructions does it operate with?
		\item Architectural: what components, agents, editors and services does it consist of?
	\end{textitemize}
\end{frame}

\begin{frame}{Main classes of interface subsystem tasks}
	\topline
	\justifying
	
	The user interface of an intelligent system solves a wide range of problems that go beyond simple input and output.
	
	\begin{textitemize}
		\item Visualization of knowledge, states, and results of the system's operation.
		\item Receiving actions, requests, commands, restrictions and clarifications.
		\item Interpretation of ambiguous, incomplete, and context-dependent messages.
		\item Support for dialogue, navigation, learning and explanation.
		\item Coordination of human interaction with internal agents of the system.
	\end{textitemize}
\end{frame}

\begin{frame}{The task of identifying and\\ interpreting the input}
	\topline
	\justifying
	
	An intelligent interface must recognize not just the shape of the input signal, but its intended meaning.
	
	\begin{textitemize}
		\item Recognition of text, speech, gestures, touches, and object manipulation.
		\item Selection of entities, relationships, intentions and task parameters.
		\item Disambiguation based on context and interaction history.
		\item Identifying errors, omissions, and conflicts in a user's message.
	\end{textitemize}
\end{frame}

\begin{frame}{The task of forming an\\ external representation}
	\topline
	\justifying
	
	The system must not only find a solution, but also present it in a form suitable for perception and further action.
	
	\begin{textitemize}
		\item Select the presentation form: text, table, diagram, graph, map, warning, animation.
		\item Select the level of detail and order of information disclosure.
		\item Adaptation of the presentation to the user's role, competence and current task.
		\item Providing interactivity and the possibility of subsequent clarification.
	\end{textitemize}
\end{frame}

\begin{frame}{Dialog Control Task}
	\topline
	\justifying
	
	The interface of an intelligent system should support not a single command, but an evolving dialogue.
	
	\begin{textitemize}
		\item Support for user initiative and system initiative.
		\item Maintaining the communication scenario and context of the current session.
		\item Transitions between clarification, explanation, teaching and performing actions.
		\item Managing dialog states, interruptions, and returning to previous steps.
	\end{textitemize}
\end{frame}

\begin{frame}{Interface Adaptation Task}
	\topline
	\justifying
	
	Intelligent interfaces adapt interaction to the characteristics of the subject and the situation.
	
	\begin{textitemize}
		\item Adaptation to the user's level of training and work style.
		\item Adaptation to the device, communication channel and physical environment.
		\item Taking into account the frequency of actions and typical scenarios.
		\item Personalize the presentation, tooltips, and order of operations.
	\end{textitemize}
\end{frame}

\begin{frame}{The task of explanation and justification}
	\topline
	\justifying
	
	It is especially important for an intelligent system to be able to explain what it has understood, why it made a decision, and what constraints were taken into account.
	
	\begin{textitemize}
		\item Explanation of the result and its basis.
		\item Display of facts, rules, hypotheses and sources used.
		\item Identification of the degree of confidence and areas of uncertainty.
		\item Support for user trust, verification, and training.
	\end{textitemize}
\end{frame}

\begin{frame}{The problem of coordination with other systems}
	\topline
	\justifying
	
	The interface of an intelligent system often works not only with humans, but also with other systems, services, and agents.
	
	\begin{textitemize}
		\item Machine-to-machine exchange of messages and knowledge.
		\item Agreement on formats, protocols, and semantics of interaction.
		\item Translation between external and internal representations.
		\item Managing distributed services and agents.
	\end{textitemize}
\end{frame}

\begin{frame}{Problem of interaction\\ with the physical environment}
	\topline
	\justifying
	
	If we are talking about a robotic or cyber-physical intelligent system, the interface subsystem must also work with the surrounding physical environment.
	
	\begin{textitemize}
		\item Receiving signals from sensors and environmental events.
		\item Formation of control actions on actuators.
		\item Taking into account spatial, temporal and situational context.
		\item Integration of human-machine and machine-environment interaction.
	\end{textitemize}
\end{frame}

\begin{frame}{Multi-level classification of interfaces}
	\topline
	\justifying
	
	It is convenient to classify intelligent interfaces according to several independent criteria at once.
	
	\begin{textitemize}
		\item By interaction recipient.
		\item By modality of communication.
		\item According to the degree of intelligence and adaptability.
		\item By the method of presentation and editing of external structures.
		\item On the architectural approach to development.
	\end{textitemize}
\end{frame}

\begin{frame}{Classification by\\ interaction recipient}
	\topline
	\justifying
	
	From the point of view of the external world, interfaces of intelligent systems can be divided according to the main interaction partner.
	
	\begin{textitemize}
		\item User interfaces: human --- system.
		\item Intersystem interfaces: system --- system.
		\item Sensory-actuator interfaces: system --- physical environment.
		\item Mixed interfaces: simultaneous coordination of all three circuits.
	\end{textitemize}
\end{frame}

\begin{frame}{Classification by communication modality}
	\topline
	\justifying
	
	Based on the form of external messages, several classes of interfaces can be distinguished.
	
	\begin{textitemize}
		\item Text and symbolic.
		\item Graphic and visual-spatial.
		\item Speech and auditory.
		\item Gestural, tactile, sensory.
		\item Multimodal, combining several channels.
	\end{textitemize}
\end{frame}

\begin{frame}{Classification by level of intelligence}
	\topline
	\justifying
	
	Conventionally, several levels of development of interface tools can be distinguished.
	
	\begin{textitemize}
		\item Traditional interfaces: fixed scenarios and controls.
		\item Adaptive interfaces: customization for the user and context.
		\item Intelligent interfaces: semantic interpretation, intent modeling, explanation.
		\item Self-evolving interfaces: accumulation of experience and knowledge-based reconfiguration.
	\end{textitemize}
\end{frame}

\begin{frame}{Traditional, adaptive, and intelligent interfaces}
	\topline
	\justifying
	
	This line of classification is particularly useful for understanding the evolution of interface tools.
	
	\begin{textitemize}
		\item The traditional interface relies on predefined forms and commands.
		\item The adaptive interface takes into account the user profile, device, preferences, and task context.
		\item The intelligent interface additionally uses knowledge and semantic interpretation.
		\item In practice, modern systems often combine all three levels.
	\end{textitemize}
\end{frame}

\begin{frame}{WIMP, Post-WIMP, and Intelligent Interfaces}
	\topline
	\justifying
	
	For computer systems, it is useful to consider the evolution of interaction paradigms.
	
	\begin{textitemize}
		\item WIMP: windows, icons, menu, pointer.
		\item Post-WIMP: gestures, touch input, direct manipulation, spatial scenes.
		\item Intelligent interfaces: semantic and dialog mechanisms on top of traditional and post-WIMP tools.
		\item Advanced systems are usually multimodal and are not limited to a single paradigm.
	\end{textitemize}
\end{frame}

\begin{frame}{Classification by the degree of specialization of the language of communication}
	\topline
	\justifying
	
	An important criterion is how universal the external language of communication is or how focused it is on a specific subject domain.
	
	\begin{textitemize}
		\item General-purpose external languages are suitable for a wide range of tasks and knowledge structures.
		\item Specialized languages are focused on specific types of objects, operations, and scenarios.
		\item Hybrids are possible between them: a universal core plus domain extensions.
	\end{textitemize}
\end{frame}

\begin{frame}{Classification by the role of external structures}
	\topline
	\justifying
	
	External information structures may differ in their function in interaction.
	
	\begin{textitemize}
		\item Representation constructs: visualization of knowledge, states, results.
		\item Assignment and editing structures: forms, graphs, diagrams, text descriptions.
		\item Navigation structures: menus, maps, routes, tables of contents, links.
		\item Explanatory constructions: comments, tracings, argumentation schemes.
	\end{textitemize}
\end{frame}

\begin{frame}{Classification by\\ construction architecture}
	\topline
	\justifying
	
	Approaches to interface development differ in the level of abstraction and the way components are described.
	
	\begin{textitemize}
		\item Manual development at the level of widgets and screen forms.
		\item Component and framework approach.
		\item Model-oriented interface design.
		\item Ontological and semantic design.
		\item Agent-oriented and event-driven architectures.
	\end{textitemize}
\end{frame}

\begin{frame}{What knowledge does\\ an intelligent interface need?}
	\topline
	\justifying
	
	Interface intelligence is impossible without the explicit or implicit use of multiple types of knowledge.
	
	\begin{textitemize}
		\item Knowledge of the subject domain.
		\item Knowledge about the user and his activities.
		\item Knowledge of interaction scenarios.
		\item Knowledge of types of external languages and constructions.
		\item Knowledge about the state of the system itself and its services.
	\end{textitemize}
\end{frame}

\begin{frame}{User model in an intelligent interface}
	\topline
	\justifying
	
	One of the central components of an intelligent interface is the user model.
	
	\begin{textitemize}
		\item Role, qualifications, preferences, limitations.
		\item History of actions and typical work scenarios.
		\item Goals, intentions, current context of the task.
		\item Possible difficulties, errors and the need for hints.
	\end{textitemize}
\end{frame}

\begin{frame}{Interaction Context Model}
	\topline
	\justifying
	
	In addition to the user model, the system must take into account the context of the situation.
	
	\begin{textitemize}
		\item Current state of the subject domain.
		\item Activity stage and active subtask.
		\item Device, communication channel, environment of use.
		\item Temporal, spatial and organizational constraints.
	\end{textitemize}
\end{frame}

\begin{frame}{Adaptation and customization\\interface}
	\topline
	\justifying
	
	Contemporary work on intelligent interfaces emphasizes the importance of adapting and personalizing user tools.
	
	\begin{textitemize}
		\item Individualization of the composition of displayed components.
		\item Changing the level of detail and sequence of actions.
		\item Reconfigure input methods, suggestions, and notifications.
		\item Formation of personalized interaction trajectories.
	\end{textitemize}
\end{frame}

\begin{frame}{User Interface Components\\ as Knowledge Objects}
	\topline
	\justifying
	
	In the ontological and model-oriented approach, user interface elements are considered as formally described entities.
	
	\begin{textitemize}
		\item Presentation components: text, images, maps, graphs, messages.
		\item Containers: windows, panels, tabs, lists, tables, trees.
		\item Interactive components: input fields, radio buttons, buttons, selectors.
		\item View manipulation components: scrolling, zooming, resizing.
	\end{textitemize}
\end{frame}

\begin{frame}{User Interface Actions}
	\topline
	\justifying
	
	For an intelligent system, it is important not only to store interface components, but also to formally describe user actions.
	
	\begin{textitemize}
		\item Mouse actions: hover, click, double-click, drag, gestures.
		\item Keyboard actions: text input, service keys, combinations.
		\item Touch actions: taps, multi-touch gestures, swipes.
		\item Speech, pen, and material actions in advanced interfaces.
	\end{textitemize}
\end{frame}

\begin{frame}{Why is formalization of interface actions important?}
	\topline
	\justifying
	
	Formal description of actions allows us to link user behavior with task semantics and internal system agents.
	
	\begin{textitemize}
		\item You can set rules for handling events at the knowledge level.
		\item Simplifies reuse and automatic interface generation.
		\item It becomes possible to analyze user scenarios and errors.
		\item Ensures consistency between the visual layer and internal processes.
	\end{textitemize}
\end{frame}

\begin{frame}{Interface and internal system agents}
	\topline
	\justifying
	
	In an intelligent system, the user interface is closely linked to internal agents that interpret events and initiate actions.
	
	\begin{textitemize}
		\item Interface events trigger computational, logical, and service procedures.
		\item Internal agents can update the view, provide explanations and clarifications.
		\item The observable form of work of intelligent agents is organized through the interface.
		\item This is especially important for explainable and dialog systems.
	\end{textitemize}
\end{frame}

\begin{frame}{Approach 1:\\ Manual interface design}
	\topline
	\justifying
	
	The simplest and historically earliest approach is direct programming of screens, forms, and event handlers.
	
	\begin{textitemize}
		\item Gives complete control over the implementation details.
		\item Works well for local and relatively stable applications.
		\item Does not scale well with increasing number of scenarios and modalities.
		\item Makes it difficult to reuse interface knowledge.
	\end{textitemize}
\end{frame}

\begin{frame}{Approach 2:\\ Component and framework}
	\topline
	\justifying
	
	Here the interface is assembled from ready-made components and interaction patterns.
	
	\begin{textitemize}
		\item The speed of development is increased.
		\item Ensures consistency of elements and behavior.
		\item Cross-platform support is simplified.
		\item However, the semantics of the subject domain often remain outside the interface description itself.
	\end{textitemize}
\end{frame}

\begin{frame}{Approach 3:\\ Model-Driven Development}
	\topline
	\justifying
	
	model-driven User Interface Development views the interface through a system of interconnected models.
	
	\begin{textitemize}
		\item Domain model.
		\item Model of user tasks and scenarios.
		\item Abstract interface model.
		\item Model of a specific interface for the selected platform.
		\item Generate a final implementation from a set of models.
	\end{textitemize}
\end{frame}

\begin{frame}{Strengths\\ of the model-driven approach}
	\topline
	\justifying
	
	This approach is important for complex intelligent systems where multiple representations and platforms need to be coordinated.
	
	\begin{textitemize}
		\item Increases description reuse.
		\item Simplifies the portability of the interface to different devices and channels.
		\item Makes explicit the relationships between tasks, data, and presentation.
		\item Facilitates partial automation of design and generation.
	\end{textitemize}
\end{frame}

\begin{frame}{Limitations of the model-driven approach}
	\topline
	\justifying
	
	Despite their advantages, task models and widgets alone are often insufficient for intelligent interaction.
	
	\begin{textitemize}
		\item It is difficult to express the rich semantics of the subject domain.
		\item It is difficult to take into account implicit knowledge and context dependencies.
		\item Means of reconciling user, context and explanation models are needed.
		\item Therefore, the role of ontological and semantic description increases.
	\end{textitemize}
\end{frame}

\begin{frame}{Approach 4:\\ Ontological interface design}
	\topline
	\justifying
	
	The ontological approach assumes that the interface and its components are described in terms of the formal ontology of the subject domain and the ontology of the interaction itself.
	
	\begin{textitemize}
		\item Component, action, view, and relationship classes are explicitly defined.
		\item The interface is associated with knowledge semantics, not just with code.
		\item A basis for automatic generation and consistency checking is provided.
		\item The expansion of the interface is simplified by expanding the ontology.
	\end{textitemize}
\end{frame}

\begin{frame}{What does the ontological approach provide?}
	\topline
	\justifying
	
	Ontological description is especially useful for intelligent systems, where the interface should not be a decorative shell, but a semantically loaded subsystem.
	
	\begin{textitemize}
		\item Coordination of the interface with the knowledge base and domain models.
		\item Automation of construction and reconfiguration of interface components.
		\item Increasing the reusability of solutions.
		\item Possibility of logical control, verification and explanation.
	\end{textitemize}
\end{frame}

\begin{frame}{User Interface Ontology}
	\topline
	\justifying
	
	The user interface ontology captures the types of components, actions, views, messages, and rules for coordinating them.
	
	\begin{textitemize}
		\item Describes output components, containers, and interactive elements.
		\item User actions and system reactions are formalized.
		\item Relationships of composition, compatibility, and nesting are defined.
		\item The connections with tasks, roles and objects of the subject domain are clarified.
	\end{textitemize}
\end{frame}

\begin{frame}{Ontological model of the graphical interface}
	\topline
	\justifying
	
	In ontological modeling of a graphical interface, the graphical entities themselves are considered as elements of the design subject domain.
	
	\begin{textitemize}
		\item Basic graphic components and their properties are highlighted.
		\item The rules for the composition of screen structures are defined.
		\item Correspondences are established between semantic entities and visual forms.
		\item This creates the basis for intelligent editors and interface generators.
	\end{textitemize}
\end{frame}

\begin{frame}{Rules and Limitations in Interface Design}
	\topline
	\justifying
	
	In an intelligent design environment, an interface can be evaluated and built not only based on templates, but also based on rules.
	
	\begin{textitemize}
		\item Component compatibility rules.
		\item Nesting restrictions and display methods.
		\item Rules for choosing the form of presentation for different types of knowledge.
		\item Rules for assessing the quality and appropriateness of an interface solution.
	\end{textitemize}
\end{frame}

\begin{frame}{Approach 5: Agent-Oriented Interface Architecture}
	\topline
	\justifying
	
	The interface can be built as a set of interacting agents, each of which is responsible for its own function.
	
	\begin{textitemize}
		\item Agent for recognition and interpretation of input messages.
		\item Agent for selecting the presentation form.
		\item Dialogue and user support agent.
		\item Adaptation and personalization agent.
		\item Agent for explaining decisions and recommendations.
	\end{textitemize}
\end{frame}

\begin{frame}{Approach 6: Event-Driven Architecture}
	\topline
	\justifying
	
	In this architecture, the interface is viewed as a system for processing external and internal events.
	
	\begin{textitemize}
		\item User and environment events are processed through a single processing loop.
		\item For different classes of events, rules of interpretation and response are specified.
		\item Conveniently integrate the interface with internal agents and services.
		\item The approach fits well with the ontological description of events and actions.
	\end{textitemize}
\end{frame}

\begin{frame}{Approach 7: Knowledge-based approach}
	\topline
	\justifying
	
	In intelligent systems, the interface is increasingly built around knowledge rather than around individual forms and screens.
	
	\begin{textitemize}
		\item The object of knowledge and its semantic connections become central.
		\item The presentation adapts to the type of knowledge and the user's task.
		\item Editing is carried out not only at the text level, but also at the level of knowledge structures.
		\item This is especially important for semantic and ontological systems.
	\end{textitemize}
\end{frame}

\begin{frame}{Approach 8: Human-Centered and Ergonomic}
	\topline
	\justifying
	
	Even in intelligent systems, interface design cannot be reduced to formal models alone.
	
	\begin{textitemize}
		\item Human cognitive capabilities and limitations are taken into account.
		\item Typical errors, attention load, and fatigue are analyzed.
		\item Clear scenarios, feedback and learning modes are designed.
		\item The intelligence of the interface should enhance convenience, not complicate work.
	\end{textitemize}
\end{frame}

\begin{frame}{External languages for communication with users}
	\topline
	\justifying
	
	The external language of communication is a system of external information structures and rules for their use, through which the user interacts with the intelligent system.
	
	\begin{textitemize}
		\item These can be text, graphic, tabular, speech and mixed forms.
		\item The language must be understandable to the user and interpretable by the system.
		\item The richer the semantics of a language, the wider the possibilities for intellectual interaction.
	\end{textitemize}
\end{frame}

\begin{frame}{Why separate external languages?}
	\topline
	\justifying
	
	The external language is not just a screen design, but the basis for meaningful communication between the user and the system.
	
	\begin{textitemize}
		\item The language specifies acceptable forms of questions, descriptions, and commands.
		\item Language determines what knowledge can be explicitly represented and edited.
		\item Language defines navigational tools, links, relationships, and contexts.
		\item The choice of language determines the power and accessibility of the interface.
	\end{textitemize}
\end{frame}

\begin{frame}{Universal External Languages}
	\topline
	\justifying
	
	Universal external languages are oriented towards the widest possible range of objects and interaction situations.
	
	\begin{textitemize}
		\item They are suitable for describing a variety of entities, relationships and constructions.
		\item Provide a unified representation of knowledge across various subject domains.
		\item Convenient as a base layer of a platform or intelligent environment.
		\item Often require additional means of simplification for the end user.
	\end{textitemize}
\end{frame}

\begin{frame}{Specialized External Languages}
	\topline
	\justifying
	
	Specialized languages are designed for a specific subject domain, class of tasks, or professional group.
	
	\begin{textitemize}
		\item They are closer to the natural ways of thinking of an expert in a given field.
		\item They allow us to hide the excessive generality of the universal language.
		\item Improves the speed of input and accuracy of interpretation of typical constructions.
		\item But they are less transferable to other areas and require support for alignment with the general semantic core.
	\end{textitemize}
\end{frame}

\begin{frame}{General-Purpose and Specialized Languages: A Comparison}
	\topline
	\justifying
	
	Typically, an effective intelligent system combines both types of external languages.
	
	\begin{textitemize}
		\item A universal language provides commonality, portability, and a common semantic foundation.
		\item Specialized language provides compactness and professional expressiveness.
		\item There must be rules of translation and coordination between them.
		\item In a good architecture, specialized languages are built as superstructures over general-purpose models.
	\end{textitemize}
\end{frame}

\begin{frame}{Examples of external languages by presentation form}
	\topline
	\justifying
	
	From a practical point of view, it is useful to distinguish several classes of external languages.
	
	\begin{textitemize}
		\item Natural-language and quasi-natural.
		\item Graph, diagram, schematic.
		\item Tabular and formulaic.
		\item Command and scenario.
		\item Multimodal languages that combine text, graphics, speech, and gestures.
	\end{textitemize}
\end{frame}

\begin{frame}{Natural Language Interface}
	\topline
	\justifying
	
	Natural language interface is one of the most prominent areas of intelligent interfaces.
	
	\begin{textitemize}
		\item The user formulates queries and messages in natural language.
		\item The system must perform syntactic, semantic and pragmatic analysis.
		\item Resolution of homonymy, synonymy, and contextual references is important.
		\item It is necessary to generate clear answers and clarifying questions.
	\end{textitemize}
\end{frame}

\begin{frame}{Stages of message processing in a natural language interface}
	\topline
	\justifying
	
	Processing a user message typically involves several layers of analysis and synthesis.
	
	\begin{textitemize}
		\item Analysis of message form: tokens, morphology, syntax.
		\item Understanding meaning: entities, relationships, intentions, context.
		\item Correlation with the internal knowledge model and the task.
		\item Synthesis of an answer, explanation, or clarifying question.
	\end{textitemize}
\end{frame}

\begin{frame}{Graph and Scheme Languages}
	\topline
	\justifying
	
	In many intelligent systems, knowledge and tasks are conveniently expressed not by linear text, but by structural visual constructions.
	
	\begin{textitemize}
		\item Semantic networks, graphs, conceptual schemes.
		\item Process, route, and dependency diagrams.
		\item Cause-and-effect and argumentation schemes.
		\item Such languages are especially useful when working with complex relationships and structures.
	\end{textitemize}
\end{frame}

\begin{frame}{Tabular and formulaic external languages}
	\topline
	\justifying
	
	For a number of tasks, compact tabular and formula forms are the most effective.
	
	\begin{textitemize}
		\item Tables of classification, comparison, conditions and parameters.
		\item Formula constructions for technical, economic and scientific fields.
		\item Combination of tables with semantic comments and checks.
		\item Such languages are convenient for precise input and editing of specialized knowledge.
	\end{textitemize}
\end{frame}

\begin{frame}{External Information Structure Editors}
	\topline
	\justifying
	
	Editors of external information structures are special tools for creating, modifying, and coordinating those structures through which the system interacts with the user or another system.
	
	\begin{textitemize}
		\item Editors for text, forms, tables, graphs, diagrams, maps, and scripts.
		\item They act not simply as graphical tools, but as part of an intelligent interface.
		\item A good editor must understand the type of construction being edited and control its correctness.
	\end{textitemize}
\end{frame}

\begin{frame}{What functions should such editors perform?}
	\topline
	\justifying
	
	An external design editor in an intelligent system usually performs more functions than a regular graphical or text editor.
	
	\begin{textitemize}
		\item Creation and modification of constructs in an acceptable language.
		\item Checking syntactic, structural and partially semantic correctness.
		\item Hints, auto-completion, navigation through related entities.
		\item Translation between specialized and universal representation forms.
	\end{textitemize}
\end{frame}

\begin{frame}{Types of External Structure Editors}
	\topline
	\justifying
	
	Depending on the language and type of tasks, several classes of such editors can be distinguished.
	
	\begin{textitemize}
		\item Text and template text editors.
		\item Graphic editors for diagrams, graphs and knowledge maps.
		\item Table editors for parameters, constraints and correspondences.
		\item Combined multimodal editors.
		\item Intelligent editors with proofreading and fragment generation.
	\end{textitemize}
\end{frame}

\begin{frame}{Intelligent Properties of Editors}
	\topline
	\justifying
	
	An editor becomes part of an intelligent interface when it does more than just draw or save shapes, but helps you work with meaning.
	
	\begin{textitemize}
		\item Suggests valid constructions based on ontology and context.
		\item Controls the constraints and type compatibility of elements.
		\item Links the edited form to internal knowledge structures.
		\item Supports explaining errors and recommending fixes.
	\end{textitemize}
\end{frame}

\begin{frame}{Connecting editors with general-purpose and specialized languages}
	\topline
	\justifying
	
	Editors are a practical place where the question of generality and specialization of language becomes particularly noticeable.
	
	\begin{textitemize}
		\item The universal editor supports a wide range of constructs.
		\item A specialized editor simplifies work in a specific subject domain.
		\item A multi-level scheme is often used: a domain editor plus translation into a universal core.
		\item This provides a balance between convenience and semantic integrity.
	\end{textitemize}
\end{frame}

\begin{frame}{Intelligent interface as a semiotic system}
	\topline
	\justifying
	
	It is useful to think of an intelligent interface as a semiotic system that links external signs, their meanings, and actions.
	
	\begin{textitemize}
		\item The syntax defines the permissible forms of constructions.
		\item Semantics connects constructions with objects and relations of the subject domain.
		\item Pragmatics determines the appropriateness of constructions in a specific interaction situation.
		\item Therefore, interface design requires working with three levels at once.
	\end{textitemize}
\end{frame}

\begin{frame}{Semantic core and visual shell}
	\topline
	\justifying
	
	In an intelligent system, it is important to distinguish between the appearance of the interface and its substantive basis.
	
	\begin{textitemize}
		\item The visual shell is responsible for the convenient form of presentation.
		\item The semantic core ensures unambiguous interpretation of constructions.
		\item Without a semantic core, adaptation and explanation quickly degenerate into superficial heuristics.
		\item Without a convenient shell, semantic power remains inaccessible to the user.
	\end{textitemize}
\end{frame}

\begin{frame}{Computer intelligent systems as a special case}
	\topline
	\justifying
	
	Although we are talking about cybernetic systems in general, many ideas receive their most advanced practical implementation in computer intelligent systems.
	
	\begin{textitemize}
		\item Graphical, web, mobile and voice interfaces.
		\item Dialogue assistants and recommendation services.
		\item Decision support systems and intelligent learning environments.
		\item Tools for working with knowledge and semantic graphs.
	\end{textitemize}
\end{frame}

\begin{frame}{Where intelligent interfaces are especially in demand}
	\topline
	\justifying
	
	They provide the greatest value where external interactions are complex, data is heterogeneous, and the cost of error is high.
	
	\begin{textitemize}
		\item Medicine and clinical decisions.
		\item Industrial and dispatch systems.
		\item Robotics and autonomous systems.
		\item Educational intellectual environments.
		\item Analytical and expert systems.
	\end{textitemize}
\end{frame}

\begin{frame}{Requirements for the quality of the user interface of an intelligent system}
	\topline
	\justifying
	
	The quality criteria here are broader than in regular interface development.
	
	\begin{textitemize}
		\item Clarity and convenience.
		\item Semantic precision and unambiguity of interpretation.
		\item Adaptability and personalization.
		\item Explainability and trust support.
		\item Consistency with internal models of knowledge and tasks.
	\end{textitemize}
\end{frame}

\begin{frame}{Typical Design Problems}
	\topline
	\justifying
	
	There are a number of persistent challenges in developing intelligent interfaces.
	
	\begin{textitemize}
		\item Gap between semantic model and visual implementation.
		\item Complexity of formalizing user intentions.
		\item Conflict between versatility and ease of use.
		\item Limited data for personalization and adaptation.
		\item Difficulties in explaining complex system outputs.
	\end{textitemize}
\end{frame}

\begin{frame}{Why You Shouldn't Limit Yourself to a GUI Approach}
	\topline
	\justifying
	
	If we consider the interface only as a graphical user layer, then a significant part of the problematic of intelligent systems is lost.
	
	\begin{textitemize}
		\item Intersystem interfaces and interfaces with the environment remain outside the scope.
		\item External languages and knowledge construct editors are underestimated.
		\item The role of ontologies, agents and semantic interpretation is not taken into account.
		\item Therefore, GUI is just one of the special cases of a more general interface subsystem.
	\end{textitemize}
\end{frame}

\begin{frame}{From screen forms to semantic interaction}
	\topline
	\justifying
	
	The general development trend is towards a transition from an interface as a set of screen forms to an interface as an intelligent intermediary between the system and the outside world.
	
	\begin{textitemize}
		\item The role of user and context models is increasing.
		\item Dialogue and multimodal means are being developed.
		\item The importance of external languages and knowledge structure editors is increasing.
		\item Design increasingly relies on models, ontologies and rules.
	\end{textitemize}
\end{frame}

\begin{frame}{Conclusion}
	\topline
	\justifying
	
	\vspace{10mm}
	
	The user interface of an intelligent system should be understood as part of a more general interface subsystem of a cybernetic system.
	
	\begin{textitemize}
		\item Intelligent interface works with meaning, context, adaptation and explanation.
		\item Its classification is multidimensional: addressee, modality, language, architecture, degree of intelligence.
		\item Key development approaches: manual, component, model-oriented, ontological, agent-based, knowledge-driven.
		\item Of particular importance are external communication languages and editors of external information structures.
	\end{textitemize}
\end{frame}


